% Part: first-order-logic
% Chapter: natural-deduction
% Section: quantifier-rules

\documentclass[../../../include/open-logic-section]{subfiles}

\begin{document}

\olfileid[pt]{fol}{ntd}{qrl}

\olsection{Regras de Quantificadores}

\subsection{Regras para $\lforall$}

\begin{defish}
\AxiomC{$!A(a)$}
\RightLabel{\Intro{\lforall}}
\UnaryInfC{$\lforall[x][\Atom{!A}{x}]$}
\DisplayProof
\hfill
\AxiomC{$\lforall[x][\Atom{!A}{x}]$}
\RightLabel{\Elim{\lforall}}
\UnaryInfC{$!A(t)$}
\DisplayProof
\end{defish}

Nas regras para~$\lforall$, $t$ é um termo fechado (um termo que não
contém nenhuma variável), e $a$~é !!a{constant} que não
ocorre na conclusão~$\lforall[x][!A(x)]$, nem em qualquer hipótese
que esteja !!{undischarged} na !!{derivation} que termina com a
premissa~$!A(a)$. Chamamos $a$ de \emph{variável própria} da
inferência \Intro{\lforall}.\footnote{Usamos o termo ``variável própria''
embora $a$ na regra acima seja uma constante. Isso tem razões
históricas.}

\subsection{Regras para $\lexists$}

\begin{defish}
\AxiomC{$\Atom{!A}{t}$}
\RightLabel{\Intro{\lexists}}
\UnaryInfC{$\lexists[x][\Atom{!A}{x}]$}
\DisplayProof
\hfill
\AxiomC{$\lexists[x][\Atom{!A}{x}]$}
\AxiomC{[$\Atom{!A}{a}$]$^n$}
\DeduceC{$!C$}
\DischargeRule{\Elim{\lexists}}{n}
\BinaryInfC{$!C$}
\DisplayProof
\end{defish}

Novamente, $t$ é um termo fechado, e $a$ é !!a{constant} que não
ocorre na premissa $\lexists[x][!A(x)]$, na conclusão~$!C$, nem em
qualquer hipótese que esteja !!{undischarged} nas !!{derivation}s que terminam
com as duas premissas (além das hipóteses $!A(a)$). Chamamos
$a$ de \emph{variável própria} da inferência \Elim{\lexists}.

A condição de que uma variável própria não ocorra nem nas premissas nem
em qualquer hipótese que esteja !!{undischarged} nas !!{derivation}s
que levam às premissas para a inferência \Intro{\lforall} ou \Elim{\lexists}
é chamada de \emph{condição de variável própria}.

\begin{explain}
Lembre-se da convenção de que, quando $!A$ é !!a{formula} com a
!!{variable}~$x$ livre, indicamos isso escrevendo~$!A(x)$. No
mesmo contexto, $!A(t)$ é então uma abreviação de~$\Subst{!A}{t}{x}$. Assim,
poderíamos também escrever a regra $\Intro\lexists$ como:
\begin{prooftree}
\AxiomC{$\Subst{!A}{t}{x}$}
\RightLabel{\Intro{\lexists}}
\UnaryInfC{$\lexists[x][!A]$}
\end{prooftree}
Note que $t$ pode já ocorrer em~$!A$, por exemplo, $!A$~poderia
ser~$\Atom{\Obj P}{t,x}$.  Assim, inferir $\lexists[x][\Atom{\Obj
P}{t,x}]$ de~$\Atom{\Obj P}{t,t}$ é uma aplicação correta
de~$\Intro\lexists$---pode-se ``substituir'' uma ou mais, e não
necessariamente todas, as ocorrências de~$t$ na premissa pela
!!{variable}~$x$ ligada. Contudo, as condições de variável própria em
$\Intro\lforall$ e~$\Elim\lexists$ exigem que a !!{constant}~$a$
não ocorra em~$!A$. Portanto, não se pode inferir corretamente
$\lforall[x][\Atom{\Obj P}{a,x}]$ de $\Atom{\Obj P}{a,a}$
usando~$\Intro\lforall$.
\end{explain}

\begin{explain}
Em \Intro{\lexists} e \Elim{\lforall} não há restrições, e
o termo~$t$ pode ser qualquer coisa, de modo que não temos de nos
preocupar com nenhuma condição. Por outro lado, nas regras \Elim{\lexists} e
\Intro{\lforall}, a condição de variável própria exige que a
!!{constant}~$a$ não ocorra em lugar algum na conclusão nem em uma
hipótese !!{undischarged}. A condição é necessária para garantir que
o sistema seja correto, isto é, que só !!{derive}s !!{sentence}s de
hipóteses !!{undischarged}s das quais elas se seguem. Sem essa
condição, o seguinte seria permitido:
\begin{prooftree}
  \AxiomC{$\lexists[x][!A(x)]$}
  \AxiomC{$\Discharge{!A(a)}{1}$}
  \RightLabel{*\Intro{\lforall}}
  \UnaryInfC{$\lforall[x][!A(x)]$}
  \RightLabel{\Elim{\lexists}}
  \BinaryInfC{$\lforall[x][!A(x)]$}
\end{prooftree}
Contudo, $\lexists[x][!A(x)] \Entails/ \lforall[x][!A(x)]$.

Como as regras de eliminação para quantificadores só permitem substituir
!!{variable}s por termos fechados, segue-se que qualquer !!{formula} que
possa ser derivada de um conjunto de !!{sentence}s é ela própria !!a{sentence}.
\end{explain}


\end{document}
